\chapter{Instalasi Perangkat Lunak dan Setup Archi MCP}
\label{app:setup-archi-mcp}

\lstset{
  basicstyle=\ttfamily\scriptsize,
  breaklines=true,
  columns=fullflexible,
  frame=lines,
  showstringspaces=false
}

Lampiran ini memandu pembaca menyiapkan lingkungan kerja untuk latihan
arsitektur enterprise: Python, Visual Studio Code (VS Code), Codex CLI, Claude
Code CLI, Archi, dan plug-in Archi MCP Server. Instruksi diverifikasi terhadap
dokumentasi resmi pada 1 Juni 2026. Karena versi perangkat lunak berubah secara
berkala, gunakan tautan resmi dan pilih rilis stabil terbaru ketika mengunduh.

\section{Hasil Akhir yang Diharapkan}

Setelah mengikuti lampiran ini, pembaca memiliki lingkungan kerja berikut:

\begin{enumerate}
  \item Python dan VS Code untuk mengelola berkas, skrip, dan folder proyek.
  \item Codex CLI dan Claude Code CLI sebagai klien AI berbasis terminal.
  \item Archi untuk membuat serta menyimpan model ArchiMate.
  \item Plug-in Archi MCP Server yang berjalan di dalam Archi.
  \item Sambungan lokal dari Claude Code atau Codex menuju model yang sedang
        terbuka di Archi melalui endpoint
        \url{http://127.0.0.1:18090/mcp}.
\end{enumerate}

\begin{table}[htbp]
\centering
\scriptsize
\caption{Perangkat lunak yang digunakan dalam lingkungan kerja.}
\label{tab:software-setup}
\begin{tabularx}{\textwidth}{@{}p{0.19\textwidth} p{0.26\textwidth} X@{}}
\hline
\textbf{Perangkat Lunak} & \textbf{Fungsi} & \textbf{Tautan Resmi} \\
\hline
Python & Menjalankan skrip bantu dan lingkungan virtual. &
\href{https://www.python.org/downloads/}{python.org/downloads} \\
Visual Studio Code & Menyunting berkas proyek dan membuka terminal terpadu. &
\href{https://code.visualstudio.com/download}{code.visualstudio.com/download} \\
Codex CLI & Klien AI OpenAI berbasis terminal. &
\href{https://developers.openai.com/codex/cli}{developers.openai.com/codex/cli} \\
Claude Code CLI & Klien AI Anthropic berbasis terminal. &
\href{https://code.claude.com/docs/en/setup}{code.claude.com/docs/en/setup} \\
Archi & Aplikasi lintas platform untuk pemodelan ArchiMate. &
\href{https://www.archimatetool.com/download/}{archimatetool.com/download} \\
Archi MCP Server & Plug-in yang mengekspos model Archi melalui protokol MCP. &
\href{https://github.com/fanievh/archi-mcp-server}{github.com/fanievh/archi-mcp-server} \\
\hline
\end{tabularx}
\end{table}

\section{Persiapan Umum}

Siapkan koneksi internet, hak untuk memasang aplikasi pada komputer, dan satu
folder kerja. Codex CLI memerlukan akun ChatGPT yang mendukung Codex atau API
key. Claude Code memerlukan akun Anthropic yang sesuai atau penyedia API yang
didukung. Jalankan perintah instalasi pada terminal biasa; gunakan hak
administrator hanya ketika sistem operasi memintanya.

Repositori Archi MCP Server mensyaratkan Archi 5.7 atau lebih baru dan Java 21
atau lebih baru. Paket Archi terbaru sudah menyertakan \emph{Java Runtime}
yang sesuai. Pada 1 Juni 2026, rilis Archi stabil terbaru adalah 5.9.0 dan
menyertakan Java Runtime 21. Jika memakai distribusi Archi khusus tanpa runtime
terbundel, pasang Java 21 dari
\url{https://adoptium.net/temurin/releases/?version=21}.

\section{Instalasi pada Windows 10 atau Windows 11}

\subsection{Python}

Buka \url{https://www.python.org/downloads/}, unduh \emph{Python install
manager} untuk Windows, lalu jalankan pemasangnya. Alternatifnya, gunakan
PowerShell:

\begin{lstlisting}
winget install 9NQ7512CXL7T -e --accept-package-agreements --disable-interactivity
py install default
py --version
\end{lstlisting}

Perintah \texttt{py install default} memasang versi Python stabil bawaan yang
direkomendasikan oleh pengelola Python.

\subsection{Visual Studio Code}

Buka \url{https://code.visualstudio.com/download}, pilih \emph{Windows User
Installer} sesuai arsitektur komputer (umumnya x64), jalankan berkas pemasang,
dan terima opsi untuk menambahkan perintah \texttt{code} ke \emph{PATH}. Tutup
dan buka kembali terminal, lalu periksa:

\begin{lstlisting}
code --version
\end{lstlisting}

\subsection{Codex CLI}

Jalankan PowerShell dan gunakan pemasang resmi Codex CLI:

\begin{lstlisting}
powershell -ExecutionPolicy ByPass -c "irm https://chatgpt.com/codex/install.ps1 | iex"
codex --version
\end{lstlisting}

Jalankan perintah \texttt{codex}. Pilih masuk dengan akun ChatGPT atau API key,
lalu ikuti petunjuk pada layar.

\subsection{Claude Code CLI}

Jalankan PowerShell:

\begin{lstlisting}
irm https://claude.ai/install.ps1 | iex
claude --version
claude doctor
\end{lstlisting}

Jalankan \texttt{claude} dan ikuti alur masuk pada peramban. Git for Windows
bersifat opsional tetapi direkomendasikan agar Claude Code dapat memakai Git
Bash. Unduh melalui
\href{https://git-scm.com/download/win}{halaman resmi Git for Windows}.

\subsection{Archi}

Buka \url{https://www.archimatetool.com/download/}, pilih \emph{Windows
Installer}, lalu jalankan pemasang. Halaman resmi Archi meminta pengguna
menghapus versi lama sebelum memasang versi baru. Jalankan Archi dari menu
Start setelah instalasi selesai.

\section{Instalasi pada Ubuntu Linux}

\subsection{Python dan Utilitas Dasar}

Buka terminal, lalu jalankan:

\begin{lstlisting}
sudo apt update
sudo apt install -y python3 python3-venv python3-pip curl
python3 --version
\end{lstlisting}

\subsection{Visual Studio Code}

Halaman unduhan VS Code menyediakan paket \texttt{.deb} untuk Ubuntu. Cara
ringkas lain yang tercantum pada halaman resmi adalah memakai Snap:

\begin{lstlisting}
sudo snap install --classic code
code --version
\end{lstlisting}

\subsection{Codex CLI dan Claude Code CLI}

Gunakan pemasang native resmi:

\begin{lstlisting}
curl -fsSL https://chatgpt.com/codex/install.sh | sh
codex --version

curl -fsSL https://claude.ai/install.sh | bash
claude --version
claude doctor
\end{lstlisting}

Jalankan \texttt{codex} dan \texttt{claude} satu per satu untuk menyelesaikan
autentikasi.

\subsection{Archi}

Unduh arsip Linux 64-bit dari
\href{https://www.archimatetool.com/download/}{archimatetool.com/download}.
Ekstrak arsip ke folder yang mudah ditemukan. Jalankan program \texttt{Archi}
di dalam folder hasil ekstraksi. Contoh:

\begin{lstlisting}
tar -xzf Archi-Linux64-*.tgz
cd Archi
./Archi
\end{lstlisting}

\section{Instalasi pada macOS}

\subsection{Python dan Visual Studio Code}

Unduh pemasang Python terbaru dari \url{https://www.python.org/downloads/}.
Paket Python resmi untuk macOS berbentuk \texttt{.pkg} universal yang
mendukung Mac Apple Silicon dan Intel. Jalankan pemasang, lalu periksa:

\begin{lstlisting}
python3 --version
\end{lstlisting}

Unduh VS Code dari \url{https://code.visualstudio.com/download}. Pilih berkas
\texttt{.dmg} yang sesuai, lalu pindahkan aplikasi ke folder
\texttt{Applications}. Gunakan paket universal bila jenis prosesor belum
diketahui.

\subsection{Codex CLI dan Claude Code CLI}

Buka Terminal dan jalankan:

\begin{lstlisting}
curl -fsSL https://chatgpt.com/codex/install.sh | sh
codex --version

curl -fsSL https://claude.ai/install.sh | bash
claude --version
claude doctor
\end{lstlisting}

Jalankan \texttt{codex} dan \texttt{claude} untuk menyelesaikan autentikasi.

\subsection{Archi}

Buka \url{https://www.archimatetool.com/download/}. Pilih \emph{macOS Silicon}
untuk Apple Silicon atau \emph{macOS Intel} untuk Mac Intel. Buka berkas
\texttt{.dmg}, lalu pindahkan aplikasi Archi ke folder \texttt{Applications}.

\section{Memasang Plug-in Archi MCP Server}

Langkah berikut berlaku sama pada Windows, Ubuntu Linux, dan macOS.

\begin{enumerate}
  \item Buka halaman rilis
        \url{https://github.com/fanievh/archi-mcp-server/releases}.
  \item Unduh aset terbaru dengan ekstensi \texttt{.archiplugin}. Pada 1 Juni
        2026, rilis terbaru adalah v1.4.0. Gunakan rilis stabil terbaru bila
        nomor versi sudah berubah.
  \item Buka Archi, lalu pilih menu \texttt{Help > Manage Plug-ins > Install
        New...}.
  \item Pilih berkas \texttt{.archiplugin} yang baru diunduh dan selesaikan
        instalasi.
  \item Tutup dan buka kembali Archi.
  \item Buat model baru atau buka berkas \texttt{.archimate} yang akan
        dikerjakan.
  \item Pilih menu \texttt{MCP Server > Start MCP Server}.
\end{enumerate}

Secara bawaan, server memakai endpoint lokal:

\begin{lstlisting}
http://127.0.0.1:18090/mcp
\end{lstlisting}

Alamat \texttt{127.0.0.1} berarti server hanya menerima sambungan dari komputer
yang sama. Pertahankan pengaturan ini untuk latihan. Jangan membuka server ke
jaringan umum tanpa memahami risiko akses dan konfigurasi TLS.

Pengaturan tambahan tersedia melalui menu \texttt{Window > Preferences > MCP
Server}. Jika port diubah, konfigurasi klien pada bagian berikut juga harus
diubah.

\section{Menghubungkan Claude Code ke Archi MCP}

Cara paling mudah untuk folder proyek adalah membuat berkas
\texttt{.mcp.json} pada akar folder:

\begin{lstlisting}
{
  "mcpServers": {
    "archi": {
      "type": "http",
      "url": "http://127.0.0.1:18090/mcp"
    }
  }
}
\end{lstlisting}

Alternatifnya, daftarkan server melalui terminal dari folder proyek:

\begin{lstlisting}
claude mcp add --transport http --scope project archi http://127.0.0.1:18090/mcp
claude mcp list
\end{lstlisting}

Jalankan \texttt{claude}. Claude Code akan meminta persetujuan sebelum memakai
server yang didefinisikan pada tingkat proyek. Di dalam sesi Claude Code,
gunakan perintah \texttt{/mcp} untuk memeriksa status sambungan.

\section{Menghubungkan Codex ke Archi MCP}

Codex memakai format konfigurasi TOML. Pilih lokasi konfigurasi yang sesuai:

\begin{itemize}
  \item konfigurasi pengguna: \path{~/.codex/config.toml};
  \item proyek tepercaya: \path{.codex/config.toml} pada akar proyek.
\end{itemize}

Isi berkas tersebut dengan konfigurasi berikut:

\begin{lstlisting}
[mcp_servers.archi]
url = "http://127.0.0.1:18090/mcp"
\end{lstlisting}

Jalankan \texttt{codex} dari folder proyek. Di dalam sesi Codex, gunakan
perintah \texttt{/mcp} untuk memastikan server \texttt{archi} aktif. Codex CLI
dan ekstensi IDE Codex memakai konfigurasi yang sama.

\section{Uji Sambungan}

Buka sebuah model di Archi dan nyalakan MCP Server. Kemudian, lakukan
pemeriksaan berikut.

\subsection{Uji Endpoint}

Pada Windows PowerShell, Ubuntu, atau macOS, jalankan:

\begin{lstlisting}
curl -i http://127.0.0.1:18090/mcp
\end{lstlisting}

Respons galat HTTP masih dapat menunjukkan bahwa server hidup. Misalnya,
\texttt{400 Bad Request} muncul karena perintah \texttt{curl} sederhana belum
mengirim \emph{handshake} MCP. Sebaliknya, \texttt{Connection refused}
menunjukkan bahwa server belum berjalan atau port tidak sesuai.

\subsection{Uji dari Klien AI}

Buka sesi Claude Code atau Codex yang baru, lalu minta:

\begin{lstlisting}
Berikan ringkasan model ArchiMate yang sedang terbuka di Archi.
\end{lstlisting}

Jika sambungan berhasil, klien akan memakai alat MCP untuk membaca model.
Lanjutkan dengan permintaan kecil, misalnya mencari semua elemen
\emph{Stakeholder} atau membuat sebuah view uji.

\section{Masalah Umum dan Perbaikannya}

\begin{table}[htbp]
\centering
\scriptsize
\caption{Masalah umum ketika menyiapkan Archi MCP.}
\label{tab:archi-mcp-troubleshooting}
\begin{tabularx}{\textwidth}{@{}p{0.27\textwidth} X@{}}
\hline
\textbf{Gejala} & \textbf{Perbaikan} \\
\hline
Menu \texttt{MCP Server} tidak muncul & Pasang kembali berkas
\texttt{.archiplugin}, lalu tutup dan buka kembali Archi. \\
\texttt{Connection refused} & Buka sebuah model di Archi, pilih
\texttt{MCP Server > Start MCP Server}, dan periksa port. \\
\texttt{MODEL\_NOT\_LOADED} & Pastikan model sudah terbuka. Hentikan dan
nyalakan kembali MCP Server bila perlu, lalu buka sesi klien AI yang baru. \\
Model berubah di Archi tetapi klien membaca keadaan lama & Buka sesi MCP baru.
Sesi klien mengikat keadaan model ketika sambungan dibuat. \\
Perubahan model hilang setelah Archi ditutup & Tekan \texttt{Ctrl+S} pada
Windows/Linux atau \texttt{Cmd+S} pada macOS. Perubahan melalui MCP berada di
memori sampai model disimpan. \\
Claude atau Codex tidak melihat server & Periksa akhiran
\texttt{/mcp}, format berkas konfigurasi, dan status melalui \texttt{/mcp} di
dalam klien. \\
Port \texttt{18090} sudah dipakai aplikasi lain & Ubah port melalui
\texttt{Window > Preferences > MCP Server}, lalu samakan URL pada semua
konfigurasi klien. \\
\hline
\end{tabularx}
\end{table}

\section{Daftar Periksa Ringkas}

\begin{enumerate}
  \item Pastikan Python, VS Code, Codex CLI, dan Claude Code CLI dapat
        dijalankan. Gunakan perintah pemeriksaan versi pada bagian instalasi.
  \item Jalankan Archi dan buka model \texttt{.archimate}.
  \item Pastikan plug-in Archi MCP Server sudah terpasang.
  \item Pilih \texttt{MCP Server > Start MCP Server}.
  \item Tambahkan konfigurasi Claude Code atau Codex.
  \item Buka sesi klien AI baru dan periksa \texttt{/mcp}.
  \item Uji pembacaan model dan simpan perubahan model secara eksplisit.
\end{enumerate}
